\documentclass[12pt]{article}
\usepackage[utf8]{inputenc}
\usepackage[T1]{fontenc}
\usepackage{hyperref}
\title{W\`e Northern Language Resources for K-12 Curriculum Development}
\author{Compiled by Claude Sonnet 5.5 (claude-sonnet-5-5) for Andr\'as Kornai}
\date{}
\begin{document}
\maketitle

\section{Introduction}
W\`e Northern (ISO 639-3 wob, Glottocode weno1238) is a Kru language of C\^ote d'Ivoire, also called Wob\'e, Ouob\'e or W\`e\`e. It is closely related to W\`e Southern (Gu\'er\'e). Only a few resources turned up in search, and none of them is school material.

\section{Primary Educational Resources}
No primers, textbooks or school programmes were found. The language is not known to be taught in schools.

\section{Linguistic and Reference Materials}
Ingeborg Egner, \emph{Pr\'ecis de grammaire wob\'e} (1989, Universit\'e d'Abidjan, in French), is the main grammar found. Its Glottolog record is at \url{https://glottolog.org/resource/reference/id/116019}.

J. Bertho (1951) wrote a comparative study of the G\'er\'e and Wob\^e dialects. Its Glottolog record is at \url{https://glottolog.org/resource/reference/id/27617}.

The CNRS REFLEX bibliography lists word-list and sociolinguistic works on Wob\'e: \url{https://reflex.cnrs.fr/Lexiques/webball/biblio.php?LGNOM=Wob\'e}.

Wikipedia's Wob\'e article gives a phonological overview: \url{https://en.wikipedia.org/wiki/Wob%C3%A9_language}. The article on the related Guere language is at \url{https://en.wikipedia.org/wiki/Guere_language}.

Omniglot has a page on Wob\'e (W\ae\ae): \url{https://omniglot.com/writing/wobe.htm}.

\section{Religious and Community Resources}
None found in the searches.

\section{Digital and Audio Resources}
No audio, video or software resources were found. No dictionary appeared in the results, although one may exist.

\section{Curriculum Development Context}
With about 160{,}000 speakers, W\`e Northern is a vigorous language. The documentation found is a grammar and comparative or sociolinguistic studies, mostly in French. A curriculum would need new orthography, primer and dictionary work with the community and Ivorian education authorities.
\end{document}